\chapter{Feshbach Molecules: a Doorway to the BEC-BCS Crossover}
\label{Chpt:FeshmolIntro}

The previous part of this thesis was dedicated to the identification and characterization of a suitable interspecies Feshbach resonance between dysprosium and potassium. Having established the interaction properties of the Dy-K mixture, we now turn to the study of the weakly bound DyK Feshbach molecules associated at the \SI{7}{G} resonance.

Feshbach molecules are of crucial relevance in the context of the \ac{BEC}-\ac{BCS} crossover \cite{Eagles1969ppw, Leggett1980xxx, Nozieres1985bci, Giorgini2008tou, Strinati2018tbb}. In a two-component Fermi gas, the positive-scattering-length side of a Feshbach resonance supports weakly bound dimers composed of one atom from each component, as already discussed in Chapter~\ref{Chpt:Feshbach}. These dimers are composite bosons and, if cooled to sufficiently low temperature, can condense into a \ac{mBEC} \cite{Jochim2003bec, Greiner2003eoa, Zwierlein2003oob}. This molecular condensate represents the \ac{BEC} limit of the crossover. On the opposite side of the resonance, where no two-body bound state exists in vacuum, pairing can still occur as a many-body effect, leading to extended Cooper pairs and to the \ac{BCS} limit \cite{Regal2004oor, Bartenstein2004cfa, Zwierlein2004cop, Chin2004oop}. The weakly-bound Feshbach molecules can be considered the starting point to access the attractive side of the \ac{BEC}-\ac{BCS} crossover.


In alkali systems, the molecular potentials are often sufficiently well characterized to guide the search for resonances and to interpret the observed spectra. In contrast, mixtures involving lanthanide atoms, such as our Dy-K system, display a much richer and denser interaction spectrum, arising from the anisotropic electronic structure and large magnetic moment of dysprosium, which makes coupled-channel calculations difficult or altogether unavailable. 

As described in Chapter~\ref{Chpt:Feshbach} and applied in Chapters~\ref{Chpt:Lowfieldpaper} and~\ref{Chpt:AddDyKRes}, the knowledge of the binding energy is crucial to reconstruct the near-threshold molecular branch and characterize the resonance beyond the information contained in atom-loss spectroscopy alone, especially in the case of narrow resonances, for which the knowledge of the differential magnetic moment between the open- and closed-channels molecular potentials is critical to the characterization of the universal regime. 

Feshbach molecules also play a central role in a different, though related, direction of ultracold-molecule research. In many heteronuclear alkali systems, they are used as an intermediate state for coherent transfer to the rovibrational ground state, typically through stimulated Raman adiabatic passage. This route, pioneered with bialkali molecules, has led to a series of remarkable achievements, including degenerate gases of NaK molecules \cite{Schindewolf2022eom} and \ac{mBEC} of NaCs \cite{Bigagli2024oob} and NaRb \cite{Shi2025bec}. These developments are motivated by dipolar many-body physics, quantum simulation, controlled chemistry, and precision measurements. In the work presented here, however, the weakly bound molecule itself is the main object of interest, and the path toward a molecular condensate begins with understanding and controlling the loss processes that limit the molecular lifetime.


A closely related system to our Dy-K mixture is the Li-Cr Fermi-Fermi mixture \cite{Ciamei2022ddf}. Like Dy-K, Li-Cr is a mass-imbalanced heteronuclear Fermi-Fermi mixture with similar long-term scientific goals. A key difference, however, is that the molecular potential of LiCr is sufficiently well characterized theoretically to allow quantitative predictions of the resonance positions and molecular binding energies \cite{Ciamei2022euc}. In a similar and parallel approach to the results reported in this part of the thesis, LiCr Feshbach molecules have been thoroughly characterized \cite{Finelli2024ula}. In the DyK case, no theoretical guidance is available, making the experimental characterization of the molecular states presented in the following chapters an essential prerequisite for any further study of the strongly interacting regime.

While the magneto-association of DyK Feshbach molecules is rather straightforward, the purification of the molecular sample proved more challenging. Exploiting the difference in magnetic moment between the molecules, potassium, and dysprosium, and working in  a weak \ac{ODT}, we engineered a Stern-Gerlach technique able to separate the three species, keeping the dimers trapped and spilling the non-associated atoms. In Chapter~\ref{Chpt:MolTrap}, we present this reliable protocol to associate and trap a pure sample of DyK molecules, and characterize the resonance parameters through two independent measurements of the molecular bound state properties.

In Chapter~\ref{Chpt:MolTrap}, we additionally identify light-induced losses as the dominant loss mechanism of the trapped molecules. The rate of such losses depends strongly on the trapping wavelength and on the density of excited molecular states accessible from the ground molecular state, which for a lanthanide-based molecule like DyK is exceptionally rich. We systematically explore  these processes, measuring loss rates in \ac{ODT}s operating on different wavelengths. Controlling and reducing the light-induced losses allowed us to observe an additional loss mechanism caused by inelastic collisions between molecules. Characterizing these two classes of losses, understanding their dependence on the nature of the molecular state and trapping conditions, and identifying strategies to mitigate them, constitute the central experimental effort of Chapter~\ref{Chpt:MolLoss}.
















